\chapter{Cyclone and Flood Health Risks}

\section*{Definition and Overview}
Cyclones and floods are natural disasters that bring a range of health risks, both during the event and in the weeks and months afterwards. Cyclones bring destructive winds, storm surge, and heavy rain, while floods contaminate water supplies and damage homes and infrastructure. Health risks include drowning and physical injury, infections from contaminated water, loss of access to medicines and health services, and mental health effects such as anxiety, stress, and grief. Preparation before a disaster, safe behaviour during it, and careful recovery afterwards significantly reduce harm. In Australia, cyclones affect the northern coastlines, and floods can occur across much of the country.

\section*{Epidemiology}
Cyclones and floods are among the most common natural disasters in Australia, particularly in Queensland, the Northern Territory, Western Australia, and northern New South Wales. Flooding is the most frequent natural disaster in Australia, and cyclones occur during the November to April season. Drowning is the leading cause of death in floods, and most flood deaths occur in vehicles. Beyond the immediate event, affected communities experience higher rates of injury, infection, and mental health problems for months afterwards, and the burden falls most heavily on people with existing illness, older adults, and those with limited resources.

\section*{Aetiology and Pathophysiology}
The health risks arise from the physical forces of the disaster and the disruption it causes.

\begin{itemize}
    \item \textbf{Drowning:} floodwater, storm surge, and driving into floodwaters
    \item \textbf{Physical injury:} flying debris, structural damage, and falls during clean-up
    \item \textbf{Contaminated water:} sewage, chemicals, and mud contaminate floodwater and drinking supplies
    \item \textbf{Infections:} wound infections, leptospirosis, and gastrointestinal illness from contaminated water
    \item \textbf{Mould:} damp homes promote mould growth, aggravating asthma and allergies
    \item \textbf{Disruption of care:} loss of power, medicines, and access to health services
    \item \textbf{Vector-borne disease:} standing water increases mosquitoes, raising the risk of Ross River and other infections
    \item \textbf{Psychological stress:} loss of homes, possessions, and livelihoods causes grief, anxiety, and depression
\end{itemize}

\section*{Clinical Features}
Health problems may appear during the event or days to weeks afterwards.

\begin{itemize}
    \item Injuries: cuts, bruises, fractures, and crush injuries
    \item Hypothermia from cold, wet conditions
    \item Gastrointestinal illness: diarrhoea, vomiting, and fever from contaminated water or food
    \item Skin and wound infections, including from mud and debris
    \item Respiratory symptoms from mould, dust, or smoke
    \item Exacerbation of asthma, heart disease, and diabetes from stress and disrupted care
    \item Anxiety, panic, insomnia, and grief
    \item Symptoms of depression and post-traumatic stress in the weeks and months after
\end{itemize}

\section*{Differential Diagnosis}
Symptoms after a disaster are assessed in the usual way, with attention to disaster-specific causes.

\begin{itemize}
    \item Diarrhoea from contaminated water versus food poisoning or other infection
    \item Breathing problems from mould, smoke, or dust versus asthma or pneumonia
    \item Wound infections versus cellulitis from other causes
    \item Anxiety and low mood as normal disaster stress versus clinical depression or PTSD
    \item Leptospirosis and other infections carried in floodwater
    \item Heat-related illness or hypothermia depending on conditions
\end{itemize}

\section*{When to Seek Medical Care}
Seek medical care for injuries, fevers, persistent diarrhoea, breathing problems, or significant distress after a disaster, and for ongoing management of chronic conditions.

\textbf{Call 000 (or local emergency services) for:}
\begin{itemize}
    \item Any emergency during the disaster: if you are trapped, in floodwater, or in danger
    \item Difficulty breathing or chest pain
    \item A serious injury or uncontrolled bleeding
    \item Confusion, drowsiness, or loss of consciousness
    \item Signs of severe infection: high fever, spreading redness, or severe pain
    \item Suicidal thoughts or immediate danger of self-harm
\end{itemize}

\section*{Diagnosis and Investigations}
Assessment follows symptoms, with awareness of disaster-related causes.

\begin{itemize}
    \item \textbf{History and examination:} exposure to floodwater, injuries, and symptoms
    \item \textbf{Blood and stool tests:} for infections such as leptospirosis and gastrointestinal illness
    \item \textbf{Wound assessment:} cleaning, examination, and tetanus review
    \item \textbf{Imaging:} X-rays for suspected fractures or chest problems
    \item \textbf{Mental health assessment:} screening for anxiety, depression, and post-traumatic stress
    \item \textbf{Chronic disease review:} monitoring of blood pressure, blood sugar, and medicines
\end{itemize}

\section*{Management}
Management spans the phases of the disaster.

\begin{itemize}
    \item \textbf{Before:} prepare an emergency kit, an evacuation plan, and copies of prescriptions
    \item \textbf{During:} follow emergency warnings, never drive into floodwater, and move to higher ground
    \item \textbf{After:} safe clean-up with protective gear, handwashing, and wound care
    \item \textbf{Water safety:} use bottled or boiled water until supplies are declared safe
    \item \textbf{Food safety:} discard food that has touched floodwater
    \item \textbf{Mould management:} dry homes thoroughly and clean mould safely
    \item \textbf{Medical care:} treatment of injuries, infections, and chronic conditions
    \item \textbf{Mental health support:} counselling, community support, and help-seeking for persistent distress
    \item \textbf{Recovery assistance:} insurance, financial help, and community recovery services
\end{itemize}

\section*{Complications}
\begin{itemize}
    \item Drowning, the leading cause of flood death
    \item Traumatic injury and death from structural damage
    \item Outbreaks of gastrointestinal and wound infections
    \item Leptospirosis and other waterborne diseases
    \item Worsening of asthma, heart disease, and mental illness
    \item Loss of access to medicines, dialysis, and other essential care
    \item Carbon monoxide poisoning from generators and heaters used indoors
    \item Long-term mental health problems, including PTSD and depression
    \item Financial hardship and displacement
\end{itemize}

\section*{Prognosis and Outcomes}
Most people survive cyclones and floods, and most physical injuries heal with treatment. The longer-term health picture depends on the extent of loss and the quality of recovery support. Communities that prepare well, act on warnings, and receive coordinated recovery assistance recover faster. Mental health effects are common but usually improve with time and support, although a minority develop longer-term difficulties. Attention to clean water, safe food, wound care, and chronic disease management prevents most disaster-related illness.

\section*{Prevention and Self-Care}
\begin{itemize}
    \item Prepare an emergency kit: water, food, medicines, torch, and important documents
    \item Know your risk, your local warnings, and your evacuation route
    \item Never drive, walk, or ride through floodwater
    \item Keep at least three days of medicines and a medicines list
    \item Protect homes with sandbags, roof maintenance, and insurance
    \item During clean-up, wear gloves, boots, and masks and wash hands frequently
    \item Use generators and heaters outdoors only, to prevent carbon monoxide poisoning
    \item Check on neighbours, especially older people and those with chronic illness
    \item Seek help early for ongoing distress; support services are free and confidential
\end{itemize}

\section*{Sources and Further Reading}
The following reputable sources provide further information for healthcare students and patients seeking reliable, current advice about cyclone and flood health risks.

\begin{itemize}
    \item Australian Government. \textit{Get Ready Queensland and emergency preparedness.} https://www.getready.qld.gov.au/
    \item Australian Institute for Disaster Resilience. \textit{Emergency planning resources.} https://knowledge.aidr.org.au/
    \item Healthdirect Australia. \textit{Staying safe during and after floods.} https://www.healthdirect.gov.au/
    \item World Health Organization. \textit{Floods and health.} https://www.who.int/
    \item Beyond Blue. \textit{Support after natural disasters.} https://www.beyondblue.org.au/
    \item Queensland Health. \textit{Health advice after floods.} https://www.health.qld.gov.au/
\end{itemize}
